\textbf{Ambiguity note.} The phrase ``points in general position'' is not defined in the excerpt. In planar combinatorial geometry it often means ``no three collinear''; sometimes, in circle-related problems, it also includes ``no four concyclic''. I treat both:
\begin{itemize}
\item \textbf{Literal (minimal) convention:} general position means no three points are collinear.
\item \textbf{Stronger convention:} no three collinear and no four points lie on a common circle.
\end{itemize}

Given $k\ge 3$, define $n_k$ to be the smallest integer $N$ (if it exists) such that every set of $N$ points in $\mathbb{R}^2$ in general position contains a subset $S$ of $k$ points with the property:
\begin{quote}
For all distinct triples $T\subset S$ with $|T|=3$, the circumradius of the circle through $T$ is different for different $T$.
\end{quote}
(Each triple is non-collinear under both conventions, so it determines a unique circumcircle and circumradius.)

The problem asks to determine $n_k$.
